\documentclass[10pt, aspectratio=169]{beamer}
\usepackage{amsmath, amssymb, amsthm}
\usepackage{xcolor}

\usetheme{Madrid}
\usecolortheme{whale}

% Define custom theorem environments if needed
\newtheorem{intuition}{Intuition}
\newtheorem{motivation}{Motivation}

\title[Phase 1: FHE \& Bootstrapping]{Phase 1: The FHE Dream and The Bootstrapping Blueprint}
\subtitle{Fully Homomorphic Encryption over the Integers}
\author{Paper Presentation}
\date{\today}

\begin{document}

\frame{\titlepage}

\begin{frame}{Outline: Phase 1}
    \tableofcontents
\end{frame}

\section{Motivation \& The Integer Paradigm}

\begin{frame}{Motivation: The Cloud Computing Conundrum}
    \begin{motivation}
    We possess robust cryptographic solutions for data \textbf{in transit} (TLS/SSL) and data \textbf{at rest} (AES). However, computing on that data traditionally requires decryption, creating a massive vulnerability when outsourcing to untrusted cloud servers.
    \end{motivation}
    \pause
    
    \vspace{0.3cm}
    \textbf{The Fully Homomorphic Encryption (FHE) Goal:}
    Allow an untrusted party to take encrypted inputs, perform arbitrary, unbounded computations on them, and return an encrypted result. 
    \pause
    
    \vspace{0.3cm}
    When decrypted, this result must match exactly what the data owner would have obtained by computing on the plaintext locally.
\end{frame}

\begin{frame}{The Quest for the Simplest Secure Scheme}
    \textit{"What is the simplest encryption scheme for which one can hope to achieve security?"}
    \pause
    
    \vspace{0.3cm}
    \textbf{The Caesar Cipher Analogy:}
    \begin{itemize}
        \item Shift 'A' by 3 and 'B' by 3. You can trivially add those shifts together.
        \item It is perfectly (and ironically) homomorphic under addition. 
        \item \textbf{Fatal Flaw:} It is completely insecure.
    \end{itemize}
    \pause

    \vspace{0.3cm}
    \textbf{The Paradigm Shift (2009 vs. 2010):}
    \begin{itemize}
        \item Craig Gentry's breakthrough 2009 FHE scheme utilized ideal lattices over polynomial rings—conceptually and mathematically heavy machinery.
        \pause
        \item \textbf{This Paper:} Asks if we can achieve the exact same FHE blueprint using only elementary addition and multiplication over the integers. 
    \end{itemize}
\end{frame}

\section{Formalizing Fully Homomorphic Encryption}

\begin{frame}{Formalizing the FHE Dream}
    To prove a scheme works, we must rigorously define it. A homomorphic public key encryption scheme $\mathcal{E}$ consists of four Probabilistic Polynomial-Time (PPT) algorithms:
    
    \pause
    $$\mathcal{E} = (\text{KeyGen}, \text{Encrypt}, \text{Decrypt}, \textbf{Evaluate})$$
    
    \vspace{0.3cm}
    \textbf{The Evaluate Algorithm:}
    Takes as input a public key $pk$, a boolean circuit $C$, and a tuple of ciphertexts $\vec{c}=\langle c_{1},...,c_{t}\rangle$ (corresponding to the $t$ input bits of $C$), and outputs a new evaluated ciphertext $c_{out}$.
\end{frame}

\begin{frame}{Correctness and Homomorphism}
    \begin{definition}[Correct Homomorphic Decryption]
    The scheme $\mathcal{E}$ is correct for a given $t$-input circuit $C$ if, for any key-pair $(sk,pk)\leftarrow \text{KeyGen}(\lambda)$, any plaintexts $m_{1},...,m_{t}\in\{0,1\}$, and any ciphertexts $c_{i}\leftarrow \text{Encrypt}(pk,m_{i})$, the following holds with overwhelming probability:
    $$\text{Decrypt}(sk,\text{Evaluate}(pk,C,\vec{c}))=C(m_{1},...,m_{t})$$
    \end{definition}
    \pause

    \vspace{0.3cm}
    \begin{definition}[Fully Homomorphic Encryption]
    The scheme $\mathcal{E}$ is homomorphic for a class $\mathcal{C}$ of circuits if it is correct for all circuits $C\in\mathcal{C}$. The scheme $\mathcal{E}$ is \textbf{fully homomorphic} if it is correct for \textit{all} boolean circuits (i.e., circuits of arbitrary depth and size).
    \end{definition}
\end{frame}

\begin{frame}{The "Trivial" Solution and The Compactness Rule}
    \textbf{Mathematical Loophole:} Can we satisfy FHE right now? 
    \pause
    
    \textit{Yes.} An untrusted server could just take the circuit $C$, concatenate it to the raw input ciphertexts, and send it back. The user decrypts the inputs, reads $C$, and computes it locally.
    \pause
    
    \vspace{0.3cm}
    \textbf{The Fix: Compactness.} To prevent the user from doing the computational work, the server must compress the answer.
    \pause

    \begin{definition}[Compactness]
    A homomorphic scheme $\mathcal{E}$ is compact if there exists a fixed polynomial bound $b(\lambda)$ such that the size of the ciphertext output by $\text{Evaluate}(pk, C, \vec{c})$ is strictly bounded by $b(\lambda)$ bits, \textbf{independently of the size or depth of $C$}.
    \end{definition}
\end{frame}

\section{The Bootstrapping Blueprint}

\begin{frame}{The Crisis of Noise}
    \begin{intuition}[The Noise Bottleneck]
    In all known FHE schemes, ciphertexts contain "noise" to ensure security. 
    \begin{itemize}
        \item Homomorphic addition adds the noise.
        \item Homomorphic multiplication \textbf{multiplies} the noise.
    \end{itemize}
    Because multiplication squares the noise, evaluating a deep circuit causes the noise to explode past the decryption boundary, permanently destroying the plaintext.
    \end{intuition}
    \pause
    
    \vspace{0.3cm}
    This limits us to a \textit{Somewhat} Homomorphic Encryption (SHE) scheme. How do we compute infinitely if noise acts as a strict depth limit?
\end{frame}

\begin{frame}{Defining Bootstrappability}
    The answer is Gentry's Bootstrapping Blueprint. First, we define a specific class of circuits.
    \pause

    \vspace{0.3cm}
    \begin{definition}[Augmented Decryption Circuits]
    Let $\mathcal{E}$ be an encryption scheme where decryption is implemented by a boolean circuit. The set of augmented decryption circuits, $D_{\mathcal{E}}$, consists of two circuits: 
    \begin{enumerate}
        \item Decrypt two ciphertexts and ADD the results.
        \item Decrypt two ciphertexts and MULTIPLY the results.
    \end{enumerate}
    \end{definition}
    \pause

    \vspace{0.3cm}
    \begin{definition}[Bootstrappable Encryption]
    Let $\mathcal{C}_{\mathcal{E}}$ be the set of permitted circuits a scheme can evaluate before noise exceeds the bounds. A scheme is \textbf{bootstrappable} if it can evaluate its own augmented decryption circuits:
    $$D_{\mathcal{E}} \subseteq \mathcal{C}_{\mathcal{E}}$$
    \end{definition}
\end{frame}

\begin{frame}{Gentry's Bootstrapping Theorem}
    \begin{theorem}[Gentry's Transformation]
    There is an efficient transformation that, given a description of a bootstrappable scheme $\mathcal{E}$, outputs a description of a \textbf{Fully Homomorphic} encryption scheme capable of evaluating circuits of any depth.
    \end{theorem}
    \pause
    
    \vspace{0.3cm}
    \begin{intuition}[How Bootstrapping "Refreshes" Ciphertexts]
    \begin{enumerate}
        \item Take a ciphertext $c$ whose noise is dangerously high.
        \item Encrypt $c$ \textit{again} under a new public key.
        \item Provide the server with an \textit{encrypted} secret key.
        \item The server homomorphically evaluates the decryption algorithm on the inner ciphertext.
    \end{enumerate}
    \pause
    \textbf{Result:} Because the decryption circuit is evaluated homomorphically, the output is a fresh ciphertext encrypting the exact same message, but with a fixed, smaller noise level determined solely by the depth of the decryption circuit itself!
    \end{intuition}
\end{frame}

\section{References}
\begin{frame}{References}
    \begin{thebibliography}{9}
        \bibitem{DGHV10}
        Marten van Dijk, Craig Gentry, Shai Halevi, and Vinod Vaikuntanathan.
        \textit{Fully Homomorphic Encryption over the Integers}. 
        Eurocrypt, 2010.
        
        \bibitem{Gentry09}
        Craig Gentry.
        \textit{Fully Homomorphic Encryption Using Ideal Lattices}.
        STOC, 2009.
        
        \bibitem{GentryThesis}
        Craig Gentry.
        \textit{A Fully Homomorphic Encryption Scheme}.
        PhD thesis, Stanford University, 2009.
    \end{thebibliography}
\end{frame}

\end{document}